\section{The explicit tree construction}\label{sec:tree-recipe}

We now specify the single tree that indexes the submission. Every operation
in this section corresponds to a transparent definition in
\code{Trees.lean}; the equations fix the actual tree, including its branch
order and its successor constructors. The construction has three parts:
an initial family of tree transformations, amplification through all finite
functional types, and promotion in which an amplified tree supplies an
iteration index. The hierarchy \(H\) from Section~\ref{sec:hierarchy}
then turns the resulting tree into a natural number.

Throughout, \(\W\) is the inductive type with constructors
\(\Z\), \(\Suc:\W\to\W\), and
\(\Lim:(\N\to\W)\to\W\). These are literal constructors.
In particular, a displayed limit includes its specified sequence of
children. The construction uses no equivalence relation identifying
different presentations, and every recursion below has an explicit
structural argument. We retain suggestive source names such as
\code{omegaPow} and \code{veblenTree} while defining their meaning entirely
through these equations.

\subsection{Finite iteration and tree-indexed iteration}

For an endofunction \(f:X\to X\) on any type \(X\), write
\[
 f^0(x)=x,\qquad f^{k+1}(x)=f(f^k(x)).
\]
This is \code{treeIter}; despite its name, its definition is polymorphic
and its iteration counter is an ordinary natural number. We will use it
both on trees and on functions between trees. Thus the superscript always
denotes composition iteration, with the underlying type determined by
the supplied function.

A separate operation, \code{stageIter}, uses a tree as its index. For
\(f:\W\to\W\), define
\begin{align}
 I_f(\Z,x)&=x,\nonumber\\
 I_f(\Suc(t),x)&=f(I_f(t,x)),\label{eq:tree-iteration}\\
 I_f(\Lim(s),x)&=\Lim\bigl(k\mapsto I_f(s(k),x)\bigr).\nonumber
\end{align}
The recursion is on the first tree argument \(t\). The starting tree
\(x\) remains a parameter, and each limit branch makes the recursive call
at its corresponding child. No comparison between \(x\) and \(t\) is
required. In particular, the later expression \(I_f(t,t)\) is a legitimate
application of an already defined structural iterator.

The finite trees and the first branching tree are
\[
 \widehat 0=\Z,\qquad \widehat{k+1}=\Suc(\widehat k),
 \qquad \omega_{\W}=\Lim(k\mapsto\widehat k).
\]
They are \code{stageNat} and \code{stageOmega}. Keeping the two iteration
operations distinct will matter when finite iteration is applied to a
transformation of functions, and the result is subsequently used inside
a tree-indexed iteration.

\subsection{The initial arithmetic-shaped operations}

Write \(a\oplus b\) for \code{stageAdd a b} and \(a\otimes b\)
for \code{stageMul a b}. Their defining equations are
\begin{align*}
 a\oplus\Z&=a,
 &a\oplus\Suc(b)&=\Suc(a\oplus b),
 &a\oplus\Lim(s)&=\Lim(k\mapsto a\oplus s(k)),\\
 a\otimes\Z&=\Z,
 &a\otimes\Suc(b)&=(a\otimes b)\oplus a,
 &a\otimes\Lim(s)&=\Lim(k\mapsto a\otimes s(k)).
\end{align*}
Both recurse on the right argument. The order in the multiplication
successor clause is part of the specification: the accumulated result
appears on the left of \(\oplus\), with the original \(a\) on its right.
No commutativity or algebraic simplification is being assumed.

Let \(p:\W\to\W\) denote \code{omegaPow}. It is defined by
\[
 p(\Z)=\Suc(\Z),\qquad
 p(\Suc(a))=p(a)\otimes\omega_{\W},\qquad
 p(\Lim(s))=\Lim(k\mapsto p(s(k))).
\]
These equations provide a concrete initial transformation for subsequent
construction. The names and notation record the motivation for the
recipe; the submitted number requires only the stated operations on
\(\W\). We do not replace any of these trees by an ordinal representative
or invoke arithmetic laws about such representatives.

\subsection{Single and common closure recipes}

For a tree transformation \(f\), first define the tree
\[
 O_f(x)=\Lim(k\mapsto f^k(x)).
\]
This is \code{omegaClose f x}. Its zeroth child is exactly \(x\).
The word ``closure'' names this construction; it does not assert that
applying \(f\) leaves the resulting tree unchanged.

Given any transformation \(c:\W\to\W\), define its enumeration
\(E_c:\W\to\W\) by
\begin{align*}
 E_c(\Z)&=c(\Z),\\
 E_c(\Suc(a))&=c\bigl(\Suc(E_c(a))\bigr),\\
 E_c(\Lim(s))&=\Lim(k\mapsto E_c(s(k))).
\end{align*}
This is \code{enumerateClosure}. The successor inside its middle clause
is deliberate. The transformation called \code{derivative f} is
\[
                         D(f)=E_{O_f}.
\]
These definitions make sense for every supplied total \(f\) and \(c\),
without hypotheses of monotonicity, continuity, or fixed-point behavior.
For example, the first two enumerated values of the derivative satisfy
\[
 D(f)(\Z)=O_f(\Z),\qquad
 D(f)(\Suc(\Z))=O_f\bigl(\Suc(O_f(\Z))\bigr).
\]
The second value therefore starts a fresh iteration sequence from the
successor of the first constructed tree. It is not obtained by merely
selecting a later child of the original sequence. This distinction fixes
where the successor constructors enter the seed recipe and prevents an
informal fixed-point interpretation from erasing part of its syntax.

We also require a specified recipe for a sequence
\(\mathbf f=(f_0,f_1,\ldots)\) of tree transformations. Its finite rounds
and accumulated orbit are
\begin{align*}
 r^{\mathbf f}_0(x)&=x,
 &r^{\mathbf f}_{m+1}(x)&=f_m(r^{\mathbf f}_m(x)),\\
 u^{\mathbf f,x}_0&=x,
 &u^{\mathbf f,x}_{m+1}&=
       r^{\mathbf f}_{m+1}(u^{\mathbf f,x}_m).
\end{align*}
These are \code{closureRound} and \code{closureOrbit}. Round \(m\)
applies \(f_0\) first and \(f_{m-1}\) last. The orbit repeatedly starts
the next, longer round from the result already obtained. For example,
its first two noninitial values are
\(u_1=f_0(x)\) and \(u_2=f_1(f_0(f_0(x)))\).
This example fixes both the order of composition and the repeated
application of earlier functions.

Finally set
\[
 J_{\mathbf f}(x)=\Lim(m\mapsto u^{\mathbf f,x}_m),
 \qquad D_*(\mathbf f)=E_{J_{\mathbf f}}.
\]
The corresponding names are \code{commonClose} and
\code{commonDerivative}. They are explicit branching-tree operations;
their use does not require a separate theorem about common fixed points.

\subsection{The indexed seed transformation}

For a base transformation \(b:\W\to\W\), define a family
\(V_b(a):\W\to\W\), structurally indexed by \(a\in\W\), through
\begin{align*}
 V_b(\Z)&=b,\\
 V_b(\Suc(a))&=D(V_b(a)),\\
 V_b(\Lim(s))&=D_*(k\mapsto V_b(s(k))).
\end{align*}
This family is \code{veblenTree base}. At a successor index the previous
transformation is passed to the single-transformation recipe; at a limit
index the whole sequence of recursively obtained transformations is
passed to the common recipe. The definitions specify actual functions,
so no choice of a canonical index or a preferred branch sequence occurs.

Define the transformation
\[
 K_b=D\bigl(a\mapsto V_b(a)(\Z)\bigr).
\]
This is \code{criticalLift base}. The tree used as the initial argument
later in the construction is
\begin{equation}\label{eq:gamma-tree}
                    \gamma=K_p(\Z),
\end{equation}
which is exactly \code{gammaTree}. The displayed \(\gamma\) is a name
for this tree. It carries no asserted identification with a named ordinal
or a named growth rate. All subsequent operations can therefore be
checked using this definition alone.

\subsection{Amplification at every finite functional type}

The next part builds transformations uniformly through the finite types
\[
                    X_0=\W,\qquad X_{r+1}=X_r\to X_r.
\]
These are the types \code{TreeLevel r}. A value at level two transforms
tree endofunctions; a value at level three transforms those transformations,
and the recursion continues through every natural level.

Define a sequence-combining operation \(L_r:(\N\to X_r)\to X_r\)
and a successor operation \(S_r:X_r\to X_r\) by
\begin{align*}
 L_0(s)&=\Lim(s),
 &S_0(x)&=\Suc(x),\\
 L_{r+1}(s)(x)&=L_r(k\mapsto s(k)(x)),
 &S_{r+1}(g)(x)&=S_r(g(x)).
\end{align*}
They are \code{levelLimit} and \code{levelSucc}. Above level zero they
act pointwise. Thus the apparent limit of functions is constructed by
lambda abstraction and the lower-level operation. No limit in a topology,
existence theorem, or extensional equality proof is needed to define it.

For each \(r\), the amplifier \(A_r\in X_{r+2}\) is
\begin{equation}\label{eq:typed-amplifier}
 A_r(f)(x)=S_r\bigl(L_r(k\mapsto f^k(x))\bigr),
 \qquad f\in X_{r+1},\quad x\in X_r.
\end{equation}
This is \code{omegaAmplifier r}. The finite iteration is well typed
because \(f\) is an endofunction on \(X_r\). The successor is applied
after the sequence-combining operation at the same level. Retaining this
successor is material to the growth proofs; deleting it would define a
different candidate. The construction does not assume that successive
type levels produce numerically increasing values.

The first higher functional level makes the typing especially visible.
For \(F:X_1\to X_1\), a tree transformation \(f\in X_1\), and a tree
\(a\), the pointwise equations give
\[
 A_1(F)(f)(a)=
       \Suc\bigl(\Lim(k\mapsto(F^k(f))(a))\bigr).
\]
Here the finite iteration acts on transformations before the resulting
transformation is applied to \(a\). In contrast, \(A_0(f)(a)\) iterates
\(f\) directly on trees. This explains why the varying type level adds
structure beyond increasing the natural iteration counter in a single
tree transformation. Both expressions nevertheless reduce to values of
the same raw tree type once their arguments have been supplied.

To turn these higher-type values back into trees, fix
\(f:\W\to\W\) and \(a\in\W\). Define seeds
\[
 \sigma^{f,a}_0=a,\qquad
 \sigma^{f,a}_1=f,\qquad
 \sigma^{f,a}_{r+2}=A_r,
\]
and descent maps \(C^{f,a}_r:X_r\to\W\) by
\[
 C^{f,a}_0(x)=x,\qquad
 C^{f,a}_{r+1}(g)=C^{f,a}_r\bigl(g(\sigma^{f,a}_r)\bigr).
\]
They are \code{levelSeed} and \code{descendLevel}. Each descent step
applies one function to a seed of exactly the required type, then
decreases the level. It does not inspect a function's implementation
or search for a suitable argument.

The resulting tree transformation is
\begin{equation}\label{eq:higher-tower}
 T_f(a)=\Suc\!\left(
    \Lim\bigl(r\mapsto C^{f,a}_{r+2}(A_r)\bigr)\right).
\end{equation}
This is \code{higherTower f a}. Each branch has a finite type level,
although the branch function admits every level. At branch zero,
descent unfolds to
\[
 C^{f,a}_2(A_0)=A_0(f)(a)
              =\Suc\bigl(\Lim(k\mapsto f^k(a))\bigr).
\]
This identity will supply a short, concrete growth argument without
requiring a comparison between different higher-type branches.

\subsection{Promotion and the final closed tree}

The promotion operation \(P\) transforms a tree endofunction into a
new tree endofunction:
\begin{equation}\label{eq:tree-promotion}
             P(f)(a)=\Suc\bigl(I_{T_f}(T_f(a),a)\bigr).
\end{equation}
It is \code{promoteTreeFunction}. Both roles of \(T_f\) are explicit:
\(T_f(a)\) supplies the tree index, while \(T_f\) itself is the
transformation being iterated. This reuse creates no recursive equation
for an unknown object. First \(T_f\) is constructed, and then the
previously defined iterator is applied to it and to a constructed tree.

Set
\begin{align}
 f_0(a)&=\Suc(K_p(a)),\label{eq:starting-function}\\
 G(a)&=\Suc\!\left(
       \Lim\bigl(k\mapsto(P^{\,k+1}(f_0))(a)\bigr)\right).
       \label{eq:promoted-function}
\end{align}
The source names are \code{startingTreeFunction} and
\code{promotedTreeFunction}. Here \(P^{\,k+1}\) is ordinary finite
iteration on the type \(X_1=\W\to\W\). The zeroth branch is
\(P(f_0)(a)\), not \(f_0(a)\): the exponent \(k+1\) is essential
to the stated definition and to the uniform growth proof that follows.
The functions in these branches need not be ordered against one another.

With this fixed \(G\), the remaining choices are
\begin{equation}\label{eq:final-tree}
 \boxed{\quad s=G(\gamma),\qquad
 \Omega_*=\Suc\bigl(I_G(s,s)\bigr),\qquad
 N=H_{\Omega_*}(43).\quad}
\end{equation}
The same constructed tree \(s\) supplies both the index and starting
argument of the final structural iteration. In the source, the last
tree definition is exactly
\begin{LeanCode}
def finalStage : Stage :=
  let seed := promotedTreeFunction gammaTree
  .succ (stageIter promotedTreeFunction seed seed)
\end{LeanCode}
The natural-number declaration is
\code{contender\_6 := hierarchy finalStage 43}.

The input \(43\) connects this construction to the reference comparison:
quoting any old term of depth at most \(42\) costs one additional depth
level. Every reflection at input \(43\) therefore includes that entire
reference family. The tree recipe itself contains no reference to the
previous contender. Its growth mechanism remains the same at other
numerical inputs; the following section proves the relevant strict
inequalities from local successor and limit behavior of \(H\).

This completes the finite specification. Natural recursion handles
finite rounds and type levels, and structural tree recursion handles
the tree arguments. Infinitely many possible branches are represented
by total functions with finite definitions. The claim is that these
precise operations produce the stated checked natural number; assigning
an ordinal meaning or an externally named asymptotic scale is a separate
question and supplies no premise of the comparison proof.

There are also two distinct uses of successor throughout the construction.
The constructor \(\Suc\) and its pointwise extensions \(S_r\) build trees
or functions returning trees. Only after the final tree is supplied to
\(H\) does a tree successor request the numerical reflection operation
\(\RF\). The tree transformations do not evaluate the numerical hierarchy
while choosing their indices. This separation keeps the complete recipe
typed and total independently of the size of the natural number obtained
from it.
